% ===== PRÉAMBULE CANONIQUE — à coller VERBATIM en tête de CHAQUE .tex =====
\documentclass[11pt,a4paper]{article}
\usepackage[utf8]{inputenc}\usepackage[T1]{fontenc}\usepackage{lmodern}
\usepackage[french]{babel}
\usepackage{amsmath,amssymb,mathtools}\usepackage{icomma}
\usepackage[version=4]{mhchem}          % \ce{...} : équations & formules
\usepackage{siunitx}\sisetup{locale=FR} % \qty{}{}, \si{}, \ang{}
\usepackage{chemfig}                    % \chemfig{...} : structures développées
\usepackage{xcolor}\usepackage{colortbl}
\usepackage{tikz}\usepackage{pgfplots}\pgfplotsset{compat=1.18}
\usepackage[most]{tcolorbox}\usepackage{enumitem}\usepackage{array,booktabs}
\usepackage[a4paper,margin=2.2cm]{geometry}
\usetikzlibrary{arrows.meta,calc,angles,quotes,positioning,patterns,backgrounds,shapes.geometric}
\definecolor{primary}{RGB}{222,49,99}\definecolor{primarydark}{RGB}{150,22,55}
\definecolor{defcol}{RGB}{0,114,178}\definecolor{methcol}{RGB}{52,92,148}
\definecolor{excol}{RGB}{0,135,90}\definecolor{attncol}{RGB}{200,40,55}
\tcbset{boxbase/.style={enhanced,breakable,boxrule=0.9pt,arc=2.5pt,
  left=3mm,right=3mm,top=2.3mm,bottom=2.3mm,fonttitle=\bfseries,coltitle=white}}
% --- boîtes TUTORIEL (noms EXACTS, ne pas renommer) ---
\newtcolorbox{cle}[1][]{boxbase,colback=defcol!6,colframe=defcol,
  title={\textsc{À retenir}\ifx\relax#1\relax\relax\else\textmd{ : \itshape #1}\fi}}
\newtcolorbox{methode}[1][]{boxbase,colback=methcol!7,colframe=methcol,sharp corners,
  title={\textsc{Méthode}\ifx\relax#1\relax\relax\else\textmd{ --- \itshape #1}\fi}}
\newtcolorbox{exemplebox}[1][]{boxbase,colback=excol!5,colframe=excol,
  title={\textsc{Exemple}\ifx\relax#1\relax\relax\else\textmd{ : \itshape #1}\fi}}
\newtcolorbox{piege}[1][]{boxbase,colback=attncol!6,colframe=attncol,
  title={\textsc{Piège}\ifx\relax#1\relax\relax\else\textmd{ : \itshape #1}\fi}}
\newtcolorbox{remarque}[1][]{boxbase,colback=black!4,colframe=black!45,
  title={\textsc{Remarque}\ifx\relax#1\relax\relax\else\textmd{ : \itshape #1}\fi}}
% --- boîtes EXERCICE / CORRIGÉ (noms EXACTS, auto-comptées) ---
\newtcolorbox[auto counter]{exo}[1][]{boxbase,colback=primary!4,colframe=primary,
  title={\textsc{Exercice}~\thetcbcounter\ifx\relax#1\relax\relax\else\textmd{ --- \itshape #1}\fi}}
\newtcolorbox[auto counter]{corrige}[1][]{boxbase,colback=excol!4,colframe=excol,
  title={\textsc{Corrigé}~\thetcbcounter\ifx\relax#1\relax\relax\else\textmd{ --- \itshape #1}\fi}}
\newcommand{\R}{\mathbb{R}}\newcommand{\N}{\mathbb{N}}\newcommand{\Z}{\mathbb{Z}}\newcommand{\ds}{\displaystyle}
% (un \usepackage{fancyhdr}+en-tête est possible ; il est ignoré côté web)
% =========================================================================
\begin{document}
\begin{corrige}[combustions et molécules oxygénées]
On règle C, puis H, puis O en dernier (en comptant l'oxygène déjà porté par le réactif).
\begin{enumerate}[label=\alph*)]
  \item \ce{2 C4H10 + 13 O2 -> 8 CO2 + 10 H2O} \; (C : $8$ ; H : $20$ ; O : $26=16+10$).
  \item \ce{C2H5OH + 3 O2 -> 2 CO2 + 3 H2O} \; l'éthanol apporte déjà $1$ O :
        O à droite $=4+3=7$, dont $1$ vient de l'alcool, d'où $3\,\ce{O2}$.
  \item \ce{C6H12O6 + 6 O2 -> 6 CO2 + 6 H2O} \; (O : $6+12=18=12+6$).
\end{enumerate}
\end{corrige}

\begin{corrige}[groupements polyatomiques conservés]
\begin{enumerate}[label=\alph*)]
  \item \ce{2 Al(OH)3 + 3 H2SO4 -> Al2(SO4)3 + 6 H2O} \; (Al : $2$ ; \ce{SO4} : $3$ ; puis H et O
        via l'eau).
  \item \ce{Ca(OH)2 + 2 HCl -> CaCl2 + 2 H2O} \; (Ca : $1$ ; Cl : $2$ ; H : $4=4$).
  \item \ce{Fe2O3 + 6 HCl -> 2 FeCl3 + 3 H2O} \; (Fe : $2$ ; Cl : $6$ ; O : $3=3$).
\end{enumerate}
\end{corrige}

\begin{corrige}[réactions de précipitation]
\begin{enumerate}[label=\alph*)]
  \item \ce{2 AgNO3 + CaCl2 -> 2 AgCl v + Ca(NO3)2} \; (Ag : $2$ ; \ce{NO3} : $2$ ; Cl : $2$).
  \item \ce{Pb(NO3)2 + 2 KI -> PbI2 v + 2 KNO3} \; (Pb : $1$ ; I : $2$ ; K : $2$ ; \ce{NO3} : $2$).
  \item \ce{2 Na3PO4 + 3 CaCl2 -> Ca3(PO4)2 v + 6 NaCl} \; (Na : $6$ ; \ce{PO4} : $2$ ; Ca : $3$ ;
        Cl : $6$).
\end{enumerate}
\end{corrige}

\begin{corrige}[azote, soufre et gros coefficients]
\begin{enumerate}[label=\alph*)]
  \item \ce{4 NH3 + 5 O2 -> 4 NO + 6 H2O} \; (N : $4$ ; H : $12$ ; O : $10=4+6$).
  \item \ce{2 C8H18 + 25 O2 -> 16 CO2 + 18 H2O} \; l'oxygène donne $\tfrac{25}{2}\ce{O2}$, on
        double tout (O : $50=32+18$).
  \item \ce{4 FeS2 + 11 O2 -> 2 Fe2O3 + 8 SO2} \; (Fe : $4$ ; S : $8$ ; O : $22=6+16$).
\end{enumerate}
\end{corrige}

\begin{corrige}[vérifier, corriger, sommer]
\begin{enumerate}[label=\alph*)]
  \item \textbf{Non.} Inventaire : Al ($1=1$) \checkmark, Cl ($3=3$) \checkmark, mais
        \textbf{hydrogène} : $3$ à gauche (\ce{3 HCl}) contre $2$ à droite (\ce{H2}). L'équation
        \emph{n'est pas} pondérée.
  \item On cherche le ppcm des H : la réaction produisant \ce{H2} (paquets de $2$) et consommant
        \ce{HCl} (paquets de $1$), on prend $6$ atomes H de chaque côté :
        \[
        \ce{2 Al + 6 HCl -> 2 AlCl3 + 3 H2}\qquad(\text{Al}:2,\ \text{Cl}:6,\ \text{H}:6=6).
        \]
  \item Somme des coefficients : $2+6+2+3=\boxed{13}$.
\end{enumerate}
\end{corrige}

\end{document}
